\section{AMC 10 10 2001}
        \begin{enumerate}
        	\item (\textit{AMC 10 10 2001}) The median of the list $n, n + 3, n + 4, n + 5, n + 6, n + 8, n + 10, n + 12, n + 15$
is $10$. What is the mean?


 $\textbf{(A) } 4 \qquad\textbf{(B) } 6 \qquad\textbf{(C) } 7 \qquad\textbf{(D) } 10 \qquad\textbf{(E) } 11$


 
 \textbf{Answer: }E



\textbf{Solution}

The median of the list is $10$, and there are $9$ numbers in the list, so the median must be the 5th number from the left, which is $n+6$.

 We substitute the median for $10$ and the equation becomes $n+6=10$.

 Subtract both sides by 6 and we get $n=4$.

 $n+n+3+n+4+n+5+n+6+n+8+n+10+n+12+n+15=9n+63$.

 The mean of those numbers is $\frac{9n+63}{9}$ which is $n+7$. 

 Substitute $n$ for $4$ and $4+7=\boxed{\textbf{(E) }11}$.

 


	\item (\textit{AMC 10 10 2001}) A number $x$ is $2$ more than the product of its reciprocal and its additive inverse. In which interval does the number lie?


 $\textbf{(A) }  -4\leq x\leq -2 \qquad\textbf{(B) } -2<x\leq 0 \qquad\textbf{(C) } 0<x\leq 2 \qquad\textbf{(D) } 2<x\leq 4 \qquad\textbf{(E) } 4<x\leq 6$


 
 \textbf{Answer: }C



\textbf{Solution}

We can write our \href{https://artofproblemsolving.com/wiki/index.php?title=Equation}{equation} as $x= \left(\frac{1}{x} \right) \cdot (-x) +2 = -1+2 = 1$.Therefore, $\boxed{\textbf{(C) }0 < x\le 2}$.

 


	\item (\textit{AMC 10 10 2001}) The sum of two numbers is $S$. Suppose $3$ is added to each number and then each of the resulting numbers is doubled. What is the sum of the final two numbers?


 $\textbf{(A) } 2S+3 \qquad\textbf{(B) } 3S+2 \qquad\textbf{(C) } 3S+6 \qquad\textbf{(D) } 2S+6 \qquad\textbf{(E) } 2S+12$


 
 \textbf{Answer: }E



\textbf{Solution}

Suppose the two numbers are $a$ and $b$, with $a+b=S$.Then the desired sum is$2(a+3)+2(b+3)=2(a+b)+12=2S +12$, which is answer $\boxed{\textbf{(E)}}$.

 


	\item (\textit{AMC 10 10 2001}) What is the maximum number for the possible points of intersection of a circle and a triangle?


 $\textbf{(A) } 2 \qquad\textbf{(B) } 3 \qquad\textbf{(C) } 4 \qquad\textbf{(D) } 5 \qquad\textbf{(E) } 6$


 
 \textbf{Answer: }E



\textbf{Solution 1}


                    \begin{figure}[H]
                        \centering
                        \includegraphics[width=0.5\textwidth]{Solution/AMC_10_10_2001_4_Solution_Figure_1.PNG}
                    \end{figure}
                    

 We can draw a circle and a triangle, such that each side is tangent to the circle. This means that each side would intersect the circle at one point.

 You would then have $3$ points, but what if the circle was bigger? Then, each side would intersect the circle at 2 points. 

 Therefore, $2 \times 3 = \boxed{\textbf{(E) }6}$.

 



\textbf{Solution 2}

We know that the maximum amount of points that a circle and a line segment can intersect is $2$. Therefore, because there are $3$ line segments in a triangle, the maximum amount of points of intersection is $2 \times 3 = \boxed{\textbf{(E) }6}$.

 


	\item (\textit{AMC 10 10 2001}) How many of the twelve pentominoes pictured below have at least one line of reflectional symmetry?


 
                    \begin{figure}[H]
                        \centering
                        \includegraphics[width=0.5\textwidth]{Problem/AMC_10_10_2001_5_Problem_Figure_1.png}
                    \end{figure}
                    


 $\textbf{(A) } 3 \qquad\textbf{(B) } 4 \qquad\textbf{(C) } 5 \qquad\textbf{(D) } 6 \qquad\textbf{(E) } 7$


 
 \textbf{Answer: }D



\textbf{Solution}


                    \begin{figure}[H]
                        \centering
                        \includegraphics[width=0.5\textwidth]{Solution/AMC_10_10_2001_5_Solution_Figure_1.gif}
                    \end{figure}
                    

 The ones with lines over the shapes have at least one line of symmetry. Counting the number of shapes that have line(s) on them,we find $\boxed{\textbf{(D)}\ 6}$ pentominoes.

 


	\item (\textit{AMC 10 10 2001}) Let $P(n)$ and $S(n)$ denote the product and the sum, respectively, of the digits
of the integer $n$. For example, $P(23) = 6$ and $S(23) = 5$. Suppose $N$ is a
two-digit number such that $N = P(N)+S(N)$. What is the units digit of $N$?


 $\textbf{(A) } 2 \qquad\textbf{(B) } 3 \qquad\textbf{(C) } 6 \qquad\textbf{(D) } 8 \qquad\textbf{(E) } 9$


 
 \textbf{Answer: }E



\textbf{Solution}

Denote $a$ and $b$ as the tens and units digit of $N$, respectively. Then $N = 10a+b$. It follows that $10a+b=ab+a+b$, which implies that $9a=ab$. Since $a\neq0$, $b=9$. So the units digit of $N$ is $\boxed{(\textbf{E})\ 9}$.

 


	\item (\textit{AMC 10 10 2001}) When the decimal point of a certain positive decimal number is moved four
places to the right, the new number is four times the reciprocal of the original
number. What is the original number?


 $\textbf{(A) } 0.0002 \qquad\textbf{(B) } 0.002 \qquad\textbf{(C) } 0.02 \qquad\textbf{(D) } 0.2 \qquad\textbf{(E) } 2$


 
 \textbf{Answer: }C



\textbf{Solution 1}

If $x$ is the number, then moving the decimal point four places to the right is the same as multiplying $x$ by $10000$. This gives us:  
 \[10000x=4\cdot\frac{1}{x} \implies x^2=\frac{4}{10000}\] Since 
 \[x>0\implies x=\frac{2}{100}=\boxed{\textbf{(C)}\ 0.02}\]

 



\textbf{Solution 2}

Alternatively, we could try each solution and see if it fits the problems given statements.

 After testing, we find that $\boxed{\textbf{(C)}\ 0.02}$ is the correct answer.

 


	\item (\textit{AMC 10 10 2001}) Wanda, Darren, Beatrice, and Chi are tutors in the school math lab. Their
schedule is as follows: Darren works every third school day, Wanda works
every fourth school day, Beatrice works every sixth school day, and Chi works
every seventh school day. Today they are all working in the math lab. In how
many school days from today will they next be together tutoring in the lab?


 $\textbf{(A) } 42 \qquad\textbf{(B) } 84 \qquad\textbf{(C) } 126 \qquad\textbf{(D) } 178 \qquad\textbf{(E) } 252$


 
 \textbf{Answer: }B



\textbf{Solution}

By translating the words in the problem into the language of mathematics, the problem is telling us to find the least common multiple of the four numbers given. 

 $\textrm{LCM}(3, 4, 6, 7) = \textrm{LCM}(3, 2^2, 2 \cdot 3, 7) = 2^2 \cdot 3 \cdot 7 = 84$

 So the answer is $\boxed{\textbf{(B) } 84}$.

 


	\item (\textit{AMC 10 10 2001}) The state income tax where Kristin lives is levied at the rate of $p\%$ of the first $\textdollar 28000$ of annual income plus $(p + 2)\%$ of any amount above $\textdollar 28000$. Kristin noticed that the state income tax she paid amounted to $(p + 0.25)\%$ of her annual income. What was her annual income?


 $\textbf{(A) } \textdollar 28,000\qquad\textbf{(B) } \textdollar 32,000\qquad\textbf{(C) }\textdollar 35,000\qquad\textbf{(D) } \textdollar 42,000\qquad\textbf{(E) } \textdollar 56,000$


 
 \textbf{Answer: }B



\textbf{Solution 1}

Let $A$, $T$ be Kristin's annual income and the income tax total, respectively. Notice that
 \begin{align*} T &= p\%\cdot28000 + (p + 2)\%\cdot(A - 28000) \\ &= [p\%\cdot28000 + p\%\cdot(A - 28000)] + 2\%\cdot(A - 28000) \\ &= p\%\cdot A + 2\%\cdot(A - 28000) \end{align*}
We are also given that 
 \[T = (p + 0.25)\%\cdot A = p\%\cdot A + 0.25\%\cdot A\]Thus, 
 \[p\%\cdot A + 2\%\cdot(A - 28000) = p\%\cdot A + 0.25\%\cdot A\] 
 \[2\%\cdot(A - 28000) = 0.25\%\cdot A\]Solve for $A$ to obtain $A = \boxed{\textbf{(B) }\textdollar 32000}$.

 Note: You may also assume p = 0 for easy calculation.

 


	\item (\textit{AMC 10 10 2001}) If $x$, $y$, and $z$ are positive with $xy = 24$, $xz = 48$, and $yz = 72$, then $x + y + z$ is


 $\textbf{(A) }18\qquad\textbf{(B) }19\qquad\textbf{(C) }20\qquad\textbf{(D) }22\qquad\textbf{(E) }24$


 
 \textbf{Answer: }D



\textbf{Solution 1}

The first two equations in the problem are $xy=24$ and $xz=48$. Since $xyz \ne 0$, we have $\frac{xy}{xz}=\frac{24}{48} \implies 2y=z$. We can substitute $z = 2y$ into the third equation $yz = 72$ to obtain $2y^2=72 \implies y=6$ and $2y=z=12$. We replace $y$ into the first equation to obtain $x=4$.

 Since we know every variable's value, we can substitute them in to find $x+y+z = 4+6+12 = \boxed{\textbf{(D) }22}$.

 



\textbf{Solution 2}

These equations are symmetric, and furthermore, they use multiplication. This makes us think to multiply them all. This gives $(xyz)^2 = (xy)(yz)(xz) = (24)(48)(72) = (24 \times 12)^2 \implies xyz = 288$. We divide $xyz = 288$ by each of the given equations, which yields $x = 4$, $y = 6$, and $z = 12$. The desired sum is $4+6+12 = 22$, so the answer is $\boxed{\textbf{(D) } 22}$.

 



\textbf{Solution 3(strategic guess and check)}

Seeing the equations, we notice that they are all multiples of 12. Trying in factors of 12, we find that $x = 4$, $y = 6$, and $z = 12$ work. $4 + 6+ 12 = \boxed{\textbf{(D) } 22}$

 



\textbf{Solution 4}

From the first equation, we can get $y=\frac{24}{x}$. (Note that $x$ cannot be $0$, because then $xy=0$, which is not the case). Substituting this into the third equation gives $\frac{24z}{x}=72$, so $24z=72x$ meaning that $z=3x$. Substituting our equation for $z$ into the 2nd equation gives $3x^2=48$, so $x=4$ (note that $x$ is positive), which means that $z=12$, and $y=\frac{24}{x}=6$. This means that $x+y+z=4+6+12=\boxed{\textbf{(D) } 22}$

 


	\item (\textit{AMC 10 10 2001}) Consider the dark square in an array of unit squares, part of which is shown. The first ring of squares around this center square contains $8$ unit squares. The second ring contains $16$ unit squares. If we continue this process, the number of unit squares in the $100^\text{th}$ ring is


 
                    \begin{figure}[H]
                        \centering
                        \includegraphics[width=0.5\textwidth]{Problem/AMC_10_10_2001_11_Problem_Figure_1.png}
                    \end{figure}
                    


 $\textbf{(A)}\ 396 \qquad \textbf{(B)}\ 404 \qquad \textbf{(C)}\ 800 \qquad \textbf{(D)}\ 10,\!000 \qquad \textbf{(E)}\ 10,\!404$


 
 \textbf{Answer: }C



\textbf{Solution 1}

We can partition the $n^\text{th}$ ring into $4$ rectangles: two containing $2n+1$ unit squares and two containing $2n-1$ unit squares.

 There are $2(2n+1)+2(2n-1)=4n+2+4n-2=8n$ unit squares in the $n^\text{th}$ ring.

 Thus, the $100^\text{th}$ ring has $8 \times 100 = \boxed{\textbf{(C) }800}$ unit squares.

 



\textbf{Solution 2}

We can make the $n^\text{th}$ ring by removing a square of side length $2n-1$ from a square of side length $2n+1$. 

 This ring contains $(2n+1)^2-(2n-1)^2=(4n^2+4n+1)-(4n^2-4n+1)=8n$ unit squares.

 Thus, the $100^\text{th}$ ring has $8 \times 100 = \boxed{\textbf{(C)}\ 800}$ unit squares.

 



\textbf{Solution 3 (Less Rigorous)}

Notice that the first ring around the center square contains $8$ unit squares, the second ring contains $16$ unit squares, the third contains $24$ unit squares, and so on. The number of squares in the $n^\text{th}$ ring is determined by the expression $8 \times n$. Thus, the number of unit squares in the $100^\text{th}$ ring is equal to $8 \times 100$, which equals $\boxed{\textbf{(C) }800}$ unit squares.

 -Darth\_Cadet

 


	\item (\textit{AMC 10 10 2001}) Suppose that $n$ is the product of three consecutive integers and that $n$ is divisible by $7$. Which of the following is not necessarily a divisor of $n$?


 $\textbf{(A)}\ 6 \qquad \textbf{(B)}\ 14 \qquad \textbf{(C)}\ 21 \qquad \textbf{(D)}\ 28 \qquad \textbf{(E)}\ 42$


 
 \textbf{Answer: }D



\textbf{Solution 1}

Whenever $n$ is the product of three consecutive integers, $n$ is divisible by $3!$, meaning it is divisible by $6$.

 It also mentions that it is divisible by $7$, so the number is definitely divisible by all the factors of $42$.

 In our answer choices, the one that is not a factor of $42$ is $\boxed{\textbf{(D)}\ 28}$.

 



\textbf{Solution 2}

We can look for counterexamples. For example, letting $n = 13 \cdot 14 \cdot 15$, we see that $n$ is not divisible by 28, so $\boxed{\textbf{(D) }28}$ is our answer.

 



\textbf{Solution 3(elimination)}

No matter what 3 integers you choose, one of them has to be even, so since $14 = 7 \cdot 2$, and it has 7 and 2 as a divisor, answer B is out. Now, if it wasn't divisible by 3, it could be A or C($21 = 7 \cdot 3$,and $6 = 2 \cdot 3$)m so it must be divisible by 3. Therefore, it is either D or E. Since we eliminated 6, if it was E, it would be not divisible by 6($42 = 7 \cdot 6$), but it is not, so the answer is $\boxed{\textbf{(D)}\ 28}$.

 



\textbf{Solution 4}

Let $n = a \cdot (a + 1) \cdot (a + 2)$. Minimizing $a$ so that 7 is a factor, we see that when $a = 5$, $n = 5 \cdot 6 \cdot 7$. Taking the prime factorization of this, we see that you cannot make 28 from $2 \cdot 3 \cdot 5 \cdot 7$. Therefore, the answer is $\boxed{\textbf{(D)}\ 28}$.

 


	\item (\textit{AMC 10 10 2001}) A telephone number has the form $ABC - DEF - GHIJ$, where each letter represents a different digit. The digits in each part of the numbers are in decreasing order; that is, $A > B > C$, $D > E > F$, and $G > H > I > J$. Furthermore, $D$, $E$, and $F$ are consecutive even digits; $G$, $H$, $I$, and $J$ are consecutive odd digits; and $A + B + C = 9$. Find $A$.


 $\textbf{(A)} \ 4 \qquad \textbf{(B)} \ 5 \qquad \textbf{(C)} \ 6 \qquad \textbf{(D)} \ 7 \qquad \textbf{(E)} \ 8$


 
 \textbf{Answer: }E



\textbf{Solution}

We start by noting that there are $10$ letters, meaning there are $10$ digits in total. Listing them all out, we have $0, 1, 2, 3, 4, 5, 6, 7, 8, 9$. Clearly, the most restrictive condition is the consecutive odd digits, so we create casework based on that.

 Case 1: $G$, $H$, $I$, and $J$ are $7$, $5$, $3$, and $1$ respectively. 

 A cursory glance allows us to deduce that the smallest possible sum of $A + B + C$ is $11$ when $D$, $E$, and $F$ are $8$, $6$, and $4$ respectively, so this is out of the question.

 Case 2: $G$, $H$, $I$, and $J$ are $9$, $7$, $5$, and $3$ respectively. 

 A cursory glance allows us to deduce the answer. Clearly, when $D$, $E$, and $F$ are $6$, $4$, and $2$ respectively, $A + B + C$ is $9$ when $A$, $B$, and $C$ are $8$, $1$, and $0$ respectively, giving us a final answer of $\boxed{\textbf{(E)}\ 8}$

 



\textbf{Solution 2}

The ten letters must all correspond to ten distinct digits, so every digit is used in the telephone number.

 We note that $1$, $3$, $5$, $7$, $9$ are the odd numbers, and possible sequences for $GHIJ$ are either $1$, $3$, $5$, $7$ or $3$, $5$, $7$, $9$.

 $3$, $5$, and $7$ are included in both possible sequences so that we can rule out the possibilities of $3$, $5$, and $7$ for $A$, so $A$ can only be $4$, $6$, or $8$ (Note that we can also rule out $9$ as a possibility since there is no set of $2$ unique digits that could sum with $9$ to $9$). Since every possible sequence of $DEF$ also contains $4$, we can rule out $4$ as well.

 Testing $A$=$6$ means $D$, $E$, $F$ is $0$, $2$, $4$, respectively. However, $6$ and $8$ must then be part of $ABC$, and they already sum to more than $9$. So then we can try $A$=$8$. We have two cases: Where $DEF$ includes $0$, $2$, and $4$, or when $DEF$ includes $2$, $4$, and $6$. If we say the former, $ABC$ includes $6$ as well as $8$, which, again, has a sum over $9$. Therefore, $DEF$ must include $2$, $4$, and $6$, and leaves us with $0$, $1$, and $8$ for $ABC$. This sums to 9, and since the largest digit is $8$, $\boxed{\textbf{(E)}\ 8}$ is our answer.

 


	\item (\textit{AMC 10 10 2001}) A charity sells $140$ benefit tickets for a total of $\textdollar2001$. Some tickets sell for full price (a whole dollar amount), and the rest sells for half price. How much money is raised by the full-price tickets?


 $\textbf{(A) } \textdollar 782 \qquad \textbf{(B) } \textdollar 986 \qquad \textbf{(C) } \textdollar 1158 \qquad \textbf{(D) } \textdollar 1219 \qquad \textbf{(E) }\ \textdollar 1449$


 
 \textbf{Answer: }A



\textbf{Solution 1}

Let the cost of the full price ticket be $x$, the number of full-price tickets be $A$, and the number of half-price tickets be $B$

 Let's multiply both sides of the equation that naturally follows by 2. We have

 
 \[2Ax+Bx=4002\]

 And we have $A+B=140\implies B=140-A$

 Plugging in, we get $\implies 2Ax+(140-A)(x)=4002$

 Simplifying, we get $Ax+140x=4002$

 Factoring out the $x$, we get $x(A+140)=4002\implies x=\frac{4002}{A+140}$

 We see that the fraction has to simplify to an integer (the full price is a whole dollar amount)

 Thus, $A+140$ must be a factor of 4002.

 Consider the prime factorization of $4002$: $2\times3\times23\times29$

 $A$ must be a positive integer. So, we seek a factor of $4002$ to set equal to $A+140$ so that we get an integer solution for $A$ that is less than $140$. By guess-and-check OR inspection, the appropriate factor is $174$ ($2\times3\times29$), meaning that $A$ has a value of $34$. Plug this into the above equation for $x$ to get $x = 23$.

 Therefore, the price of full tickets out of $2001$ is $23\times34=\boxed{\textbf{(A) }782}$.

 --Edits by Joseph2718 (Reason: Ease of understanding)

 



\textbf{Solution 2}

Let $f$ equal the number of full-price tickets, and let $h$ equal the number of half-price tickets. Additionally, suppose that the price of $f$ is $p$. We are trying to solve for $f \cdot p$.

 Since the total number of tickets sold is $140$, we know that 
 \[f+h=140.\] The sales from full-price tickets ($f \cdot p$) plus the sales from half-price tickets $\Big(\frac{h \cdot p}{2}$, because each hall-price ticket costs $\frac{p}{2}$ dollars$\Big)$ equals $2001.$ Then we can write 
 \[fx + \frac{hx}{2}=2001.\]

 Substituting $h=140-f$ into the second equation, we get 
 \[f \cdot p +\frac{(140-f)p}{2}=f \cdot p+\frac{140p-f\cdot p}{2}=\frac{f\cdot p+140p}{2}=2001.\]

 Multiplying by $2$ and subtracting $140p$ gives us 
 \[f\cdot p=4002-140p.\] 

 Since the problem states that $p$ is a whole number, $140p$ will be some integer multiple of $140$ that ends in a $0$. Thus, $4002-140p$ will end in a $2$. Looking at the answer choices, only $\boxed{\textbf{(A) }782}$ satisfies that condition.

 == Solution 3 This is just a Diophantine Equation!

 Let full-price sales in \# ticket be $f$.  So half-price is just $140-f$.

 Let full-price price be $p$.

 $(140-f)\frac p2+fp=2001$

 $(140-f)p+2fp=4002$

 $140p+fp=4002$

 $(140+f)p=4002$

 Factor $4002=2\times3\times23\times29$.  Note, we need a factor greater than $140$.  It is $2\times3\times29=174$, which makes $140+f=174$, $f=34$, and $p=\frac{4002}{174}=23$.

 $fp=34\times23=782=A$.

 


	\item (\textit{AMC 10 10 2001}) A street has parallel curbs $40$ feet apart. A crosswalk bounded by two parallel stripes crosses the street at an angle. The length of the curb between the stripes is $15$ feet and each stripe is $50$ feet long. Find the distance, in feet, between the stripes.


 $\textbf{(A)}\ 9 \qquad \textbf{(B)}\ 10 \qquad \textbf{(C)}\ 12 \qquad \textbf{(D)}\ 15 \qquad \textbf{(E)}\ 25$


 
 \textbf{Answer: }C



\textbf{Solution 1}

Drawing the problem out, we see we get a parallelogram with a height of $40$ and a base of $15$, giving an area of $600$.

 


 
If we look at it the other way, we see the distance between the stripes is the height and the base is $50$.

 
                    \begin{figure}[H]
                        \centering
                        \includegraphics[width=0.5\textwidth]{Solution/AMC_10_10_2001_15_Solution_Figure_1.png}
                    \end{figure}
                    

 The area is still the same, so the distance between the stripes is $600/50 = \boxed{\textbf{(C)}\ 12}$.

 



\textbf{Solution 2}

Alternatively, we could use similar triangles--the $30-40-50$ triangle (created by the length of the bordering stripe and the difference between the two curbs) is similar to the $x-y-15$ triangle, where we are trying to find $y$ (the shortest distance between the two stripes). Therefore, $y$ would have to be $\frac{40}{50} \cdot 15 = 4 \cdot 3 = \boxed{\textbf{(C)}\ 12}$.

 


	\item (\textit{AMC 10 10 2001}) The mean of three numbers is $10$ more than the least of the numbers and $15$ less than the greatest. The median of the three numbers is $5$. What is their sum?


 $\textbf{(A)} \ 5 \qquad \textbf{(B)} \ 20 \qquad \textbf{(C)} \ 25 \qquad \textbf{(D)} \ 30 \qquad \textbf{(E)} \ 36$


 
 \textbf{Answer: }D



\textbf{Solution 1}

Let $m$ be the mean of the three numbers. Then the least of the numbers is $m-10$and the greatest is $m + 15$. The middle of the three numbers is the median, $5$. So$\dfrac{1}{3}[(m-10) + 5 + (m + 15)] = m$, which can be solved to get $m=10$.Hence, the sum of the three numbers is $3\cdot 10 = \boxed{\textbf{(D) }30}$.

 



\textbf{Solution 2}

Say the three numbers are $x$, $y$ and $z$. When we arrange them in ascending order then we can assume $y$ is in the middle therefore $y = 5$. 

 We can also assume that the smallest number is $x$ and the largest number of the three is $z$.Therefore,

 
 \[\frac{x+y+z}{3} = x + 10 = z - 15\]
 \[\frac{x+5+z}{3} = x + 10 = z - 15\]

 Taking up the first equation $\frac{x+5+z}{3} = x + 10$ and simplifying we obtain $z - 2x - 25 = 0$ doing so for the equation $\frac{x+5+z}{3} = z - 15$ we obtain the equation $x - 2z + 50 = 0$

 
 \[x - 2z + 50 = 2x - 4z + 100\] 

 when solve the above obtained equation and $z - 2x - 25 = 0$ we obtain the values of $z = 25$ and $x = 0$

 Therefore the sum of the three numbers is $25 + 5 + 0 = \boxed{\textbf{(D) }30}$

 


	\item (\textit{AMC 10 10 2001}) Which of the cones listed below can be formed from a $252^\circ$ sector of a circle of radius $10$ by aligning the two straight sides?


 
                    \begin{figure}[H]
                        \centering
                        \includegraphics[width=0.5\textwidth]{Problem/AMC_10_10_2001_17_Problem_Figure_1.png}
                    \end{figure}
                    


 $\textbf{(A) } \text{A cone with slant height of 10 and radius 6}$


 $\textbf{(B) } \text{A cone with height of 10 and radius 6}$


 $\textbf{(C) } \text{A cone with slant height of 10 and radius 7}$


 $\textbf{(D) } \text{A cone with height of 10 and radius 7}$


 $\textbf{(E) } \text{A cone with slant height of 10 and radius 8}$


 
 \textbf{Answer: }C



\textbf{Solution 1}


                    \begin{figure}[H]
                        \centering
                        \includegraphics[width=0.5\textwidth]{Solution/AMC_10_10_2001_17_Solution_Figure_1.png}
                    \end{figure}
                    

 The blue lines will be joined together to form a single blue line on the surface of the cone, so $10$ will be the slant height of the cone.

 The red line will form the circumference of the base. We can compute its length and use it to determine the radius.

 The length of the red line is $\dfrac{252^{\circ}}{360^{\circ}}\cdot 2\pi \cdot 10 = \frac{7}{10} \cdot 2\pi \cdot 10 = 14\pi$. This is the circumference of a circle with radius $\frac{14\pi}{2\pi} = 7$.

 Therefore the correct answer is $\boxed{\textbf{(C)} \text{ A cone with slant height of } 10 \text{ and radius } 7}$.

 



\textbf{Solution 2}

As in Solution 1, we deduce that the slant height of the cone is $10$.

 We now recall that the lateral surface area of a cone with base radius $r$ and slant height $l$ is $\pi rl$, while the area of the sector is $\frac{252^{\circ}}{360^{\circ}} \cdot \pi \cdot 10^2 = \frac{7}{10} \cdot \pi \cdot 10^2 = 70\pi$. We must therefore have 
 \[\pi \cdot r \cdot 10 = 70\pi \iff 10r = 70 \iff r = 7,\] so the answer is $\boxed{\textbf{(C)} \text{ A cone with slant height of } 10 \text{ and radius } 7}$.

 -mathlover1205 -formatted and improved by Sevenoptimus

 


	\item (\textit{AMC 10 10 2001}) The plane is tiled by congruent squares and congruent pentagons as indicated. The percent of the plane that is enclosed by the pentagons is closest to


 
                    \begin{figure}[H]
                        \centering
                        \includegraphics[width=0.5\textwidth]{Problem/AMC_10_10_2001_18_Problem_Figure_1.png}
                    \end{figure}
                    


 $\textbf{(A)} \ 50 \qquad \textbf{(B)} \ 52 \qquad \textbf{(C)} \ 54 \qquad \textbf{(D)} \ 56 \qquad \textbf{(E)} \ 58 \qquad$


 
 \textbf{Answer: }D



\textbf{Solution}

Consider any single tile:

 
                    \begin{figure}[H]
                        \centering
                        \includegraphics[width=0.5\textwidth]{Solution/AMC_10_10_2001_18_Solution_Figure_1.png}
                    \end{figure}
                    

 If the side of the small square is $a$, then the area of the tile is $\left(3a\right)^2 = 9a^2$, with $4a^2$ covered by squares, and therefore $9a^2-4a^2 = 5a^2$ by pentagons.Hence exactly $\frac{5}{9}$ of any tile is covered by pentagons, and therefore pentagons cover $\frac{5}{9}$ of the plane. When expressed as a percentage, this is $55.\overline{5}\%$, and the closest integer to this value is $\boxed{\textbf{(D) }56}$.

 



\textbf{Solution 2}

There are $6$ smaller square tiles on each side of the bigger square. We are double counting the four corner tiles, however, so we subtract by $4$. Additionally, there are $4$ inner square inside each hole of the two 8 looking figures inside the square. Thus, there are $4$*$2$*$2$ = $16$ inner squares. Thus, there are a total of $36$ squares. Because we are dealing with percentages, and therefore fractions and ratios, we can set the side length to anything we want as long as we maintain consistency throughout the entire figure. For simplicity sake, we set the side length of each small square to $1$

 Thus, the side length of the big square is $9$. Because the amount of squares is $36$, the total area of all the small squares is $36$. The area of the big square is $9$*$9$ = $81$. Hence, $\frac{36}{81}$ or $\frac{4}{9}$ of the entire figure is covered by squares. Thus, $\frac{5}{9}$ is covered by pentagons. Expressed as a decimal, this is $55.\overline{5}\%$. The closest percent value to this is $\boxed{\textbf{(D) }56}$.

 


	\item (\textit{AMC 10 10 2001}) Pat wants to buy four donuts from an ample supply of three types of donuts: glazed, chocolate, and powdered. How many different selections are possible?


 $\textbf{(A)}\ 6 \qquad \textbf{(B)}\ 9 \qquad \textbf{(C)}\ 12 \qquad \textbf{(D)}\ 15 \qquad \textbf{(E)}\ 18$


 
 \textbf{Answer: }D



\textbf{Solution 1}

Let's use \href{https://artofproblemsolving.com/wiki/index.php?title=Stars\_and\_bars}{stars and bars}.Let the donuts be represented by $O$s. We wish to find all possible combinations of glazed, chocolate, and powdered donuts that give us $4$ in all. The four donuts we want can be represented as $OOOO$. Notice that we can add two "dividers" to divide the group of donuts into three different kinds; the first will be glazed, second will be chocolate, and the third will be powdered. For example, $O|OO|O$ represents one glazed, two chocolate, and one powdered. We have six objects in all, and we wish to turn two into dividers, which can be done in $\binom{6}{2}=15$ ways. Our answer is hence $\boxed{\textbf{(D)}\ 15}$. Notice that this can be generalized to get the stars and bars (balls and urns) identity.

 



\textbf{Solution 2}

Similar to Solution 1, we can generalize the problem by using the stars and bars formula, yielding an answer of $\dbinom{4+3-1}{3-1}=\dbinom{6}{2}=\boxed{\textbf{(D)}\ 15}$.

 



\textbf{Solution 3}

Simple casework works here as well:Set up the following ratios:
 \[4:0:0\]
 \[3:1:0\]
 \[2:2:0\]
 \[2:1:1\]

 In three of these cases we see that there are two of the same ratios (so like two boxes would have $0$), and so if we swapped those two donuts, we would have the same case. Thus we get $\frac{3!}{2!1!}$ for those $3$ (You can also set it up and logically symmetry applies). For the other case where each ratio of donuts is different, we get the normal $\binom{4}{3}=6.$ Thus, our answer is $3 \cdot 3+6 = \boxed{\textbf{(D) }15}$.

 Solution by IronicNinja

 Edit by virjoy2001 & Aopsthedude

 


	\item (\textit{AMC 10 10 2001}) A regular octagon is formed by cutting an isosceles right triangle from each of the corners of a square with sides of length $2000$. What is the length of each side of the octagon?


 $\textbf{(A) } \frac{1}{3}(2000) \qquad \textbf{(B) } {2000(\sqrt{2}-1)} \qquad \textbf{(C) } {2000(2-\sqrt{2})} \qquad \textbf{(D) } {1000} \qquad \textbf{(E) } {1000\sqrt{2}}$


 
 \textbf{Answer: }B


	\item (\textit{AMC 10 10 2001}) A right circular cylinder with its diameter equal to its height is inscribed in a right circular cone. The cone has diameter $10$ and altitude $12$, and the axes of the cylinder and cone coincide. Find the radius of the cylinder.


 $\textbf{(A)}\ \frac{8}3\qquad\textbf{(B)}\ \frac{30}{11}\qquad\textbf{(C)}\ 3\qquad\textbf{(D)}\ \frac{25}{8}\qquad\textbf{(E)}\ \frac{7}{2}$


 
 \textbf{Answer: }B


	\item (\textit{AMC 10 10 2001}) In the magic square shown, the sums of the numbers in each row, column, and diagonal are the same. Five of these numbers are represented by $v$, $w$, $x$, $y$, and $z$. Find $y + z$.


 
                    \begin{figure}[H]
                        \centering
                        \includegraphics[width=0.5\textwidth]{Problem/AMC_10_10_2001_22_Problem_Figure_1.png}
                    \end{figure}
                    


 $\textbf{(A)}\ 43 \qquad \textbf{(B)}\ 44 \qquad \textbf{(C)}\ 45 \qquad \textbf{(D)}\ 46 \qquad \textbf{(E)}\ 47$


 
 \textbf{Answer: }D



\textbf{Solution 1}

We know that $y+z=2v$, so we could find one variable rather than two.

 $v+24+w=43+v$

 $24+w=43$

 $w=19$

 
                    \begin{figure}[H]
                        \centering
                        \includegraphics[width=0.5\textwidth]{Solution/AMC_10_10_2001_22_Solution_Figure_1.png}
                    \end{figure}
                    

 $44+x=24+x+z \implies z=20$

 
                    \begin{figure}[H]
                        \centering
                        \includegraphics[width=0.5\textwidth]{Solution/AMC_10_10_2001_22_Solution_Figure_2.png}
                    \end{figure}
                    

 The sum per row is $25+21+20=66$.

 Thus $66-18-25=66-43=v=23$. 

 Since we needed $2v$ and we know $v=23$, $23 \times 2 = \boxed{\textbf{(D)}\ 46}$.

 



\textbf{Solution 2}

$v+24+w=43+v$

 $24+w=43$

 $w=19$

 
                    \begin{figure}[H]
                        \centering
                        \includegraphics[width=0.5\textwidth]{Solution/AMC_10_10_2001_22_Solution_Figure_1.png}
                    \end{figure}
                    

 $44+x=24+x+z \implies z=20$

 
                    \begin{figure}[H]
                        \centering
                        \includegraphics[width=0.5\textwidth]{Solution/AMC_10_10_2001_22_Solution_Figure_2.png}
                    \end{figure}
                    

 The magic sum is determined by the bottom row. $25+20+21=66$.

 Solving for $y$:

 $y=66-19-21=66-40=26$.

 To find our answer, we need to find $y+z$. $y+z=20+26 = \boxed{\textbf{(D)}\ 46}$.

 



\textbf{Solution 3 (Really Easy Solution)}

A nice thing to know is that any $3$ numbers that go through the middle form an arithmetic sequence.

 Using this, we know that $x=(24+z)/2$, or $2x=24+z$ because $x$ would be the average.

 We also know that because $x$ is the average the magic sum would be $3x$, so we can also write the equation $3x-46=z$ using the bottom row.

 Solving for x in this system we get $x=22$, so now using the arithmetic sequence knowledge we find that $y=26$ and $z=20$.

 Adding these we get $\boxed{\textbf{(D)}\ 46}$.

 
-harsha12345

 



\textbf{Solution 4 (Systems of Equations)}

Create an equation for every row, column, and diagonal. Let $e$ be the sum of the rows, columns, and diagonals. 
 \[w+v+24=e\]
 \[x+y+18=e\]
 \[z+46=e\]
 \[v+43=e\]
 \[x+z+24=e\]
 \[w+y+21=e\]
 \[x+w+25=e\]
 \[x+v+21=e\].

 Notice that $z+46=e$ and $x+z+24=e$ both have $z$. Equate them and you get that $x=22$.Using that same strategy, we use $v+43=e$ instead. $w+v+24=e$ is good for our purposes. It turns out that $w=19$. Since we already know those numbers, and $x+w+25=e$, We can say that $e$ will be $66$. We are now able to solve: $x+z+24=e$, $w+y+21=e$, $x+v+21=e$, and $x+y+18=e$. Respectively, $v=23$, $w=19$, $x=22$, $y=26$, and $z=20$. We only require The sum of $y+z$, which is $26+20=46$.We get that the sum of $y$ and $z$ respectively is $\boxed{\textbf{(D)}\ 46}$

 -OofPirate

 



\textbf{Solution 5}

Using the middle row and middle column, we can get $18+x+y=24+x+z$, so $y-z=6$. Then, by using the choices, $y+z=43, 44, 45, 46,$ or $47$. After solving for $y$ and $z$, we can fill up the square to see that the only possibility is $\boxed{\textbf{(D)}\ 46}$.

 


	\item (\textit{AMC 10 10 2001}) A box contains exactly five chips, three red and two white. Chips are randomly removed one at a time without replacement until all the red chips are drawn or all the white chips are drawn. What is the probability that the last chip drawn is white?


 $\textbf{(A)}\ \frac{3}{10}\qquad\textbf{(B)}\ \frac{2}{5}\qquad\textbf{(C)}\ \frac{1}{2}\qquad\textbf{(D)}\ \frac{3}{5}\qquad\textbf{(E)}\ \frac{7}{10}$


 
 \textbf{Answer: }D



\textbf{Solution 1}

Imagine that we draw all the chips in random order, i.e., we do not stop when the last chip of a color is drawn. To draw out all the white chips first, the last chip left must be red, and all previous chips can be drawn in any order. Since there are $3$ red chips, the probability that the last chip of the five is red (and so also the probability that the last chip drawn is white) is $\boxed{\textbf{(D) } \frac {3}{5}}$.

 



\textbf{Solution 2}

Let's assume we don't stop picking until all of the chips are picked. To satisfy this condition, we have to arrange the letters: $W, W, R, R, R$ such that both $W$'s appear in the first $4$. We find the number of ways to arrange the white chips in the first $4$ and divide that by the total ways to choose all the chips. The probability of this occurring is $\dfrac{\dbinom{4}{2}}{\dbinom{5}{2}} = \boxed{\textbf{(D) }\dfrac35}$

 



\textbf{Solution 3}

The amount of ways to end with a white chip is by having $RRWW, RWW,$ and $WW$. The amount of arrangements for $RRWW$ with $W$ at the end is $3$, the number of arrangements of $RWW$ with $W$ at the end is $2$, and the number of arrangements with $WW$ is just $1$. This gives us $6$ total ways to end with white. Next, the cases to end with a red are $RWRR$, and $RRR$. $RWRR$ gives us $3$ ways and $RRR$ gives us $1$ way. So the number of ways to end with a red is $4$. Thus, our answer is simply $\frac{6}{4+6}$ = $\boxed{\textbf{(D) }\dfrac35}$

 



\textbf{Solution 4}

First, we figure out the total number of ways to arrange the $5$ chips. There are $\frac{5!}{3!2!} = 10$ ways of arranging them. Now if the last chip is white, all we need to do is to arrange the 4 chips before the last $W$, which are $WWRR$. There are $\frac{4!}{2!2!} = 6$ ways of arranging the first 4 chips. Therefore, the answer is $\frac{6}{10} = \boxed{\textbf{(D) }\dfrac35}$

 



\textbf{Solution 5}

Test a few possibilities and see that all of the possibilities with white last have probability $\frac{1}{10}$, and there are $6$ possible ways, so it is $\frac{6}{10} = \boxed{\text{(D) }\frac{3}{5}}$.

 


	\item (\textit{AMC 10 10 2001}) In trapezoid $ABCD$, $\overline{AB}$ and $\overline{CD}$ are perpendicular to $\overline{AD}$, with $AB+CD=BC$, $AB<CD$, and $AD=7$. What is $AB\cdot CD$?


 $\textbf{(A)}\ 12 \qquad \textbf{(B)}\ 12.25 \qquad \textbf{(C)}\ 12.5 \qquad \textbf{(D)}\ 12.75 \qquad \textbf{(E)}\ 13$


 
 \textbf{Answer: }B



\textbf{Solution 1}


                    \begin{figure}[H]
                        \centering
                        \includegraphics[width=0.5\textwidth]{Solution/AMC_10_10_2001_24_Solution_Figure_1.png}
                    \end{figure}
                    

 If $AB=x$ and $CD=y$, then $BC=x+y$. By the \href{https://artofproblemsolving.com/wiki/index.php?title=Pythagorean\_theorem}{Pythagorean theorem}, we have $(x+y)^2=(y-x)^2+49.$ Solving the equation, we get $4xy=49 \implies xy = \boxed{\textbf{(B)}\ 12.25}$.

 



\textbf{Solution 2}

Simpler is just drawing the trapezoid and then using what is given to solve.Draw a line parallel to $\overline{AD}$ that connects the longer side to the corner of the shorter side. Name the bottom part $x$ and top part $a$.By the Pythagorean theorem, it is obvious that $a^{2} + 49 = (2x+a)^{2}$ (the RHS is the fact the two sides added together equals that). Then, we get $a^2 + 49 = 4x^2 + 4ax + a^2$, cancel out and factor and we get $49 = 4x(x+a)$. Notice that $x(x+a)$ is what the question asks, so the answer is $\boxed{\textbf{(B)}\ 12.25}$.

 Solution by IronicNinja

 



\textbf{Solution 3 (Detailed Simplification)}

We know it is a trapezoid and that $\overline{AB}$ and $\overline{CD}$ are perpendicular to $\overline{AD}$. If they are perpendicular to $\overline{AD}$ that means this is a right-angle trapezoid (search it up if you don't know what it looks like or you can look at the trapezoid in the first solution). We know $\overline{AD}$ is $7$. We can then set the length of $\overline{AB}$ to be $x$ and the length of $\overline{DC}$ to be $y$. $\overline{BC}$ would then be $x+y$. Let's draw a straight line down from point $B$ which is perpendicular to $\overline{DC}$ and parallel to $\overline{AD}$. Let's name this line $M$. Then let's name the point at which line $M$ intersects $\overline{DC}$ point $E$. Line $M$ partitions the trapezoid into ▭ $ADEB$ and $\triangle$ $BEC$. We will use the triangle to solve for $xy$ using the Pythagorean theorem. The line segment $\overline{EC}$ would be $y-x$ because $\overline{DC}$ is $y$ and $\overline{DE}$ is $x$. $\overline{DE}$ is $x$ because it is parallel to $\overline{AB}$ and both are of equal length. Because of the Pythagorean theorem, we know that $(EC)^2+(BE)^2=(BC)^2$. Substituting the values we have we get $(y-x)^2+(7)^2=(x+y)^2$. Simplifying this we get $(y^2-2xy+x^2)+(49)=(x^2+2xy+y^2)$. Now we get rid of the $x^2$ and $y^2$ terms from both sides to get $(-2xy)+(49)=(2xy)$. Combining like terms we get $(49)=(4xy)$. Then we divide by $4$ to get $(12.25)=(xy)$. Now we know that $xy\ =\ 12.25$ which is answer choice $\boxed{\textbf{(B)}\ 12.25}$.

 Solution By: MATHCOUNTSCMS25

 Fixed $\text{\LaTeX}$ - Mliu630XYZ, palaashgang, anshulb, JoyfulSapling

 



\textbf{Solution 4 (EZ Cheez)}

Choose any value for $BC$, and then use Pythagorean theorem to get $CD - AB$, and $AB = (BC - (CD - AB))/2$). Then multiply $AB \cdot CD$. 

 For example: 

 $BC=25$. $CD - AB = \sqrt{25^2 - 7^2} = 24$. $AB = (25 - 24)/2=0.5$. $CD = 0.5 + 24 = 24.5$.$AB \cdot CD = (0.5)(24.5)= \boxed{\textbf{(B)}\ 12.25}$.

 



\textbf{Solution 5 (Quick Solve)}

Denote $AB$ to $x$, and $CD$ to $y$. This means that $BC = x+y$. This means that if we draw an altitude from $B$ $CD$ and we call the foot of that $P$, we have that $BP$ is $7$ due to symmetry. This means 
 \[7^2+(y-x)^2 = (x+y)^2\]. Simplifying gets you: 
 \[49 +y ^2-2xy+x^2 = x^2+2xy+y^2 \Longrightarrow xy = \dfrac{49}{4} \Longrightarrow xy = \boxed{\textbf{(B)}\ 12.25}\].

 -jb2015007

 


	\item (\textit{AMC 10 10 2001}) How many positive integers not exceeding $2001$ are multiples of $3$ or $4$ but not $5$?


 $\textbf{(A)}\ 768 \qquad \textbf{(B)}\ 801 \qquad \textbf{(C)}\ 934 \qquad \textbf{(D)}\ 1067 \qquad \textbf{(E)}\ 1167$


 
 \textbf{Answer: }B



\textbf{Solution 1}

Out of the numbers $1$ to $12$ four are divisible by $3$ and three by $4$, counting $12$ twice.Hence $6$ out of these $12$ numbers are multiples of $3$ or $4$. 

 The same is obviously true for the numbers $12k+1$ to $12k+12$ for any positive integer $k$.

 Hence out of the numbers $1$ to $60=5\cdot 12$ there are $5\cdot 6=30$ numbers that are divisible by $3$ or $4$.Out of these $30$, the numbers $10$, $15$, $20$, $30$, $40$, $45$ and $60$ are divisible by $5$.Therefore in the set $\{1,\dots,60\}$ there are precisely $30-6=24$ numbers that satisfy all criteria from the problem statement.

 Again, the same is obviously true for the set $\{60k+1,\dots,60k+60\}$ for any positive integer $k$.

 We have $1980/60 = 33$, hence there are $24\cdot 33 = 792$ good numbers among the numbers $1$ to $1980$. At this point we already know that the only answer that is still possible is $\boxed{\textbf{(B)}}$, as we only have $20$ numbers left.

 By examining the remaining $20$ by hand we can easily find out that exactly $9$ of them match all the criteria, giving us $792+9=\boxed{\textbf{(B) }801}$ good numbers.This is correct.

 



\textbf{Solution 2}

We can solve this problem by finding the cases where the number is divisible by $3$ or $4$, then subtract from the cases where none of those cases divide $5$. To solve the ways the numbers divide $3$ or $4$ we find the cases where a number is divisible by $3$ and $4$ as separate cases. We apply the floor function to every case to get $\left\lfloor \frac{2001}{3} \right\rfloor$, $\left\lfloor \frac{2001}{4} \right\rfloor$, and $\left\lfloor \frac{2001}{12} \right\rfloor$. The first two floor functions were for calculating the number of individual cases for $3$ and $4$. The third case was to find any overlapping numbers. The numbers were $667$, $500$, and $166$, respectively. We add the first two terms and subtract the third to get $1001$. The first case is finished.

 The second case is more or less the same, except we are applying $3$ and $4$ to $5$. We must find the cases where the first case over counts multiples of five. Utilizing the floor function again on the fractions $\left\lfloor \frac{2001}{3\cdot5} \right\rfloor$, $\left\lfloor \frac{2001}{4\cdot5} \right\rfloor$, and $\left\lfloor \frac{2001}{3\cdot4\cdot5} \right\rfloor$ yields the numbers $133$, $100$, and $33$. The first two numbers counted all the numbers that were multiples of either four with five or three with five less than $2001$. The third counted the overlapping cases, which we must subtract from the sum of the first two. We do this to reach $200$. Subtracting this number from the original $1001$ numbers procures $\boxed{\textbf{(B)}\ 801}$.

 



\textbf{Solution 3}

First find the number of such integers between 1 and 2000 (inclusive) and then add one to this result because 2001 is a multiple of $3$.

 There are $\frac45\cdot2000=1600$ numbers that are not multiples of $5$.  $\frac23\cdot\frac34\cdot1600=800$ are not multiples of $3$ or $4$, so $800$ numbers are.  $800+1=\boxed{\textbf{(B)}\ 801}$

 



\textbf{Solution 4}

Take a good-sized sample of consecutive integers; for example, the first $25$ positive integers. Determine that the numbers $3, 4, 6, 8, 9, 12, 16, 18, 21,$ and $24$ exhibit the properties given in the question. $25$ is a divisor of $2000$, so there are $\frac{10}{25}\cdot2000=800$ numbers satisfying the given conditions between $1$ and $2000$. Since $2001$ is a multiple of $3$, add $1$ to $800$ to get $800+1=\boxed{\textbf{(B)}\ 801}$.

 



\textbf{Solution 5}

By PIE, there are $1001$ numbers that are multiples of $3$ or $4$ and less than or equal to $2001$. $80\%$ of them will not be divisible by $5$, and by far the closest number to $80\%$ of $1001$ is $\boxed{\textbf{(B)}\ 801}$.

 



\textbf{Solution 5}

Similar to some of the above solutions. We can divide $2001$ by $3$ and $4$ to find the number of integers divisible by $3$ and $4$. Hence, we find that there are $667$ numbers less than $2001$ that are divisible by $3$, and $500$ numbers that are divisible by $4$. However, we will need to subtract the number of multiples of $15$ from 667 and that of $20$ from $500$, since they're also divisible by 5 which we don't want. There are $133$ + $100$ = $233$ such numbers. Note that during this process, we've subtracted the multiples of $60$ twice because they're divisible by both $15$ and $20$, so we have to add $33$ back to the tally (there are $33$ multiples of $60$ that does not exceed $2001$). Lastly, we have to subtract multiples of both $3$ AND $4$ since we only want multiples of either $3$ or $4$. This is tantamount to subtracting the number of multiples of $12$. And there are $166$ such numbers. Let's now collect our numbers and compute the total: $667$ + $500$ - $133$ - $100$ + $33$ - $166$ = $\boxed{\textbf{(B)}\ 801}$.

 



\textbf{Solution 6 (Venn Diagram)}


                    \begin{figure}[H]
                        \centering
                        \includegraphics[width=0.5\textwidth]{Solution/AMC_10_10_2001_25_Solution_Figure_1.png}
                    \end{figure}
                    

 We have the following Venn diagram (sorry for the bad asymptote). To calculate this Venn diagram we can just use PIE to find the numbers or just work our way out from the inside. We just need to find the numbers that are divisible by $3$ or $4$ but not by $5$, and looking back at the diagram we just have the answer of 
 \[401+133+267 = \boxed{\textbf{(B)}\ 801}\].

 -jb2015007

 



\textbf{Solution 6}

Similar to @above:Let the function $M_{2001}(n)$ return how many multiples of $n$ are there not exceeding $2001$. Then we have that the desired number is:
 \[M_{2001}(3)+M_{2001}(4)-M_{2001}(3\cdot 4)-M_{2001}(3 \cdot 5) - M_{2001}(4 \cdot 5)+M_{2001}(3 \cdot 4 \cdot 5)\]

 Evaluating each of these we get:
 \[667+500-166-133-100+33 = 1100-299 = 801.\]

 Thus, the answer is $\boxed{\textbf{(B)}\ 801}.$

 -ConfidentKoala4

 


 



\textbf{Solution 7}

The smallest common multiple of 3, 4, and 5 is 30. There are 12 integers between 1~30 that are multiples of 3 and 4 but not 5. 2001/30 = 66 R 21. Between 1~21 there are 9 integers that are multiples of 3 and 4 but not 5. 66 * 12 + 9 = $\boxed{\textbf{(B)}\ 801}.$

 - mathlover1205

 - AH2025 (math and number corrections)

 



\textbf{Solution 8}

Let $A$, $B$, and $C$ be the sets of multiples, not exceeding 2001, of $3$, $4$, and $5$, respectively. We desire $|(A\cup B)\backslash C|$. By PIE,\begin{align*}|(A\cup B)\backslash C| &= |A| + |B| - |A \cap B| - |A \cap C| - |B \cap C| + |A \cap B \cap C| \\ &= \lfloor \frac{2001}{3} \rfloor + \lfloor \frac{2001}{4} \rfloor - \lfloor \frac{2001}{3\cdot 4} \rfloor - \lfloor \frac{2001}{3\cdot 5} \rfloor - \lfloor \frac{2001}{4\cdot 5} \rfloor + \lfloor \frac{2001}{3\cdot 4\cdot 5} \rfloor \\&= \boxed{\textbf{(B)}\ 801}\end{align*}
~J.zz.z

 


\end{enumerate}